\section{Discussion}

\subsection{Why the Naive Prior Fails}

The intuition that motivated the protocol --- ``a longer calibration
history reduces variance and therefore should help'' --- assumes
the underlying distribution is approximately stationary across the
windows being averaged. In the simulated bi-weekly regime, that
assumption is false: at $K = 100$ a single window drains $\sim 8\%$
of the unpatched (host, CVE) backlog, and the highest-EPSS pairs
preferentially. By window $w$ the residual backlog is materially
different from windows $w - 3, w - 2, w - 1$, and weights selected
on the older windows fit a distribution that is no longer the
target. \textbf{Smoothing reduces variance but increases bias, and
under fast distributional shift the bias term dominates.}

The Paper~7 lag-1 hazard at $K = 200$ is therefore a \textit{lower
bound} on the calibration penalty, not an upper bound. Any procedure
that averages over more than one window of past calibration state
will pay a larger bias penalty than lag-1. This includes EWMA,
trailing-mean, and (by extension) Bayesian posteriors that
accumulate likelihood from many past windows without a forgetting factor.

\subsection{What Does Work}

The result by elimination points to procedures that adapt the
\textit{amount} of history to the rate of distributional change.
Three candidate families that are out of scope for this paper but
that the rejection pattern motivates:

\textbf{(i) Change-point detection.} Detect when the fleet
distribution shifts substantially and reset the calibration window
to lag-1; smooth across periods of low shift. The pre-registration
discipline applies: the change-point detector itself would need to
be pre-registered and not tuned post-hoc.

\textbf{(ii) Forgetting-factor Bayesian update.} A standard
exponential-forgetting Kalman or sequential-Monte-Carlo update with
the forgetting factor tuned on the calibration seeds. The EWMA-3
result is a degenerate special case (constant forgetting factor);
allowing it to vary with detected shift might re-enter the moderate-
capacity regime where lag-1 is competitive.

\textbf{(iii) Capacity-conditional deployment.} At high capacity,
turn the recalibrator off and accept the fixed Paper~4 weights,
incurring the Paper~6 absolute-floor collapse rather than the
deeper Paper~7 + Paper~8 calibration hazard. This is the simplest
deployable rule that follows from the present rejection pattern.

\subsection{Operational Implications}

Sharpened from Paper~7:
\begin{itemize}
\item Deploy lag-1 online recalibration only at $K \leq 100$ on the
simulated bi-weekly cadence. The Paper~7 K=200 hazard is real, and
Paper~8 shows that smoothing does not rescue it.
\item Do not deploy any multi-window-history smoother (EWMA-3,
trail3, or by extension longer-history variants without an adaptive
forgetting mechanism) on this calibration grid in any regime studied.
\item For organisations at high capacity, the deployable options
collapse to (a) fixed Paper~4 weights, (b) periodic offline
rebaselining between deployment windows, or (c) adaptive
recalibration methods outside the scope of this paper.
\end{itemize}

\subsection{Relationship to the Research Sequence}

Paper~4 established hygiene augmentation; Paper~5 added temporal
robustness; Paper~6 added capacity sensitivity; Paper~7 added a
deployable calibration substitute; this paper closes the easy door
on extending Paper~7. The cumulative arc is a five-step refinement
of an operational claim: hygiene-augmented prioritization is better
than EPSS-only across all studied regimes (per-pair); its absolute
P@50 floor is regime-dependent (Paper~6); deployable recalibration
is regime-conditional (Paper~7); and smoothing the recalibrator
does not extend the regime in which it is safe to deploy (this paper).

\subsection{What This Paper is Not}

This paper does not show that smoothing is intrinsically harmful in
all calibration regimes. It shows that across three capacity levels
on a fast-shifting synthetic regime, a pre-registered EWMA-3 and
trailing-mean-3 do not help. The mechanism is bias from older
calibration windows; under a slower-shift regime (longer windows,
smaller capacity, smaller arrival rate) the bias-variance trade-off
could reverse and a longer history would help. We do not survey that
slower-shift regime here.
